%==============================================================================
%  WEEKLY AI BRIEF -- issue #03, window 2-8 September 2026
%  Engine: tectonic (also compiles with pdflatex/latexmk)
%==============================================================================
\documentclass[11pt,a4paper]{article}

%------------------------------ packages -------------------------------------
\usepackage[T1]{fontenc}
\usepackage[utf8]{inputenc}
\usepackage{lmodern}
\usepackage[margin=2.2cm,top=2.0cm,bottom=2.3cm]{geometry}
\usepackage{microtype}
\usepackage{graphicx}
\usepackage{xcolor}
\usepackage{titlesec}
\usepackage{enumitem}
\usepackage{booktabs}
\usepackage{tabularx}
\usepackage{fancyhdr}
\usepackage{lastpage}
\usepackage{tcolorbox}
\usepackage{needspace}
\usepackage[hidelinks]{hyperref}

\tcbuselibrary{skins,breakable}

% Logos live in template/ and are shared by every issue -- never copy them into a
% dated folder. This path list resolves them from <YYYY-MM-DD>/ or from template/.
\graphicspath{{../../../template/}{../template/}{./}{template/}}

%------------------------------ palette --------------------------------------
\definecolor{brandink}{HTML}{103868}    % KAUST navy -- headings, rules
\definecolor{brandaccent}{HTML}{0B7480} % KAUST teal (darkened to 5.6:1 on white)
\definecolor{brandsoft}{HTML}{EAF3F5}   % box fills
\definecolor{brandgrey}{HTML}{5A6672}   % meta text

\hypersetup{colorlinks=true,linkcolor=brandaccent,urlcolor=brandaccent,
            citecolor=brandaccent,pdfborder={0 0 0}}

%------------------------------ metadata -------------------------------------
\newcommand{\briefnumber}{03}
\newcommand{\briefweek}{2--8 September 2026}
\newcommand{\briefdate}{8 September 2026}
\newcommand{\briefauthor}{Ali Habibullah}
\newcommand{\briefaffil}{KAUST}
\newcommand{\briefscope}{Papers, model releases, benchmarks and policy from the past seven days}

%------------------------------ section styling ------------------------------
\titleformat{\section}
  {\normalfont\sffamily\large\bfseries\color{brandink}}{\thesection.}{0.6em}{}
  [\vspace{-0.55em}{\color{brandaccent}\rule{\textwidth}{0.9pt}}]
\titlespacing*{\section}{0pt}{1.5ex plus .3ex}{0.9ex}

\titleformat{\subsection}[runin]
  {\normalfont\sffamily\bfseries\color{brandink}}{}{0pt}{}[\;---]
\titlespacing*{\subsection}{0pt}{1.25ex plus .2ex}{0.5em}

\setlist[itemize]{leftmargin=1.15em,itemsep=0.25em,topsep=0.35em,parsep=0pt}
\setlength{\parindent}{0pt}
\setlength{\parskip}{0.45em}

%------------------------------ header / footer ------------------------------
\pagestyle{fancy}
\fancyhf{}
\renewcommand{\headrulewidth}{0pt}
\renewcommand{\footrulewidth}{0.4pt}
\fancyfoot[L]{\footnotesize\color{brandgrey}Weekly AI Brief \#\briefnumber\ \textbullet\ \briefweek}
\fancyfoot[R]{\footnotesize\color{brandgrey}Page \thepage\ of \pageref{LastPage}}

%------------------------------ content macros -------------------------------
\newcommand{\tldr}[1]{%
  \begin{tcolorbox}[breakable,enhanced,colback=brandsoft,colframe=brandaccent,
      boxrule=0.8pt,left=10pt,right=10pt,top=8pt,bottom=8pt,arc=2pt,
      title={\sffamily\bfseries The week in five lines},fonttitle=\color{white}]
    \begin{itemize}[leftmargin=1.1em,itemsep=0.3em,topsep=0pt]#1\end{itemize}
  \end{tcolorbox}}

\newcommand{\entry}[4]{%
  \par\vspace{0.7em}%
  \Needspace*{5\baselineskip}%
  {\sffamily\bfseries\color{brandink}#1}\par
  {\footnotesize\color{brandgrey}#2 \textbullet\ \href{#3}{link} \textbullet\ \srcref{#4}}\par
  \vspace{0.15em}}

\newcommand{\why}[1]{\par{\small\textbf{Why it matters.} #1}\par}

\newcommand{\srcref}[1]{\textcolor{brandaccent}{[S#1]}}

\newcommand{\stat}[2]{{\sffamily\bfseries\color{brandink}#1}~{\footnotesize\color{brandgrey}#2}}

%==============================================================================
\begin{document}

%------------------------------ title block ----------------------------------
\noindent
\begin{minipage}[c]{0.48\textwidth}
  \raggedright\includegraphics[height=1.4cm]{KAUST_Academy_Logo.png}
\end{minipage}\hfill
\begin{minipage}[c]{0.48\textwidth}
  \raggedleft\includegraphics[height=1.5cm]{KAUST_Logo.png}
\end{minipage}
\vspace{0.7em}

{\sffamily
 {\color{brandgrey}\footnotesize ISSUE \#\briefnumber\ \textbullet\ \briefdate}\par
 \vspace{0.15em}
 {\LARGE\bfseries\color{brandink}Weekly AI Brief}\par
 \vspace{0.25em}
 {\large\color{brandaccent}\briefweek}\par
 \vspace{0.35em}
 {\footnotesize\color{brandgrey}\briefscope\\
  Prepared by \briefauthor, \briefaffil}\par}
\vspace{0.3em}
{\color{brandink}\rule{\textwidth}{1.4pt}}
\vspace{0.6em}

%------------------------------ 1. executive summary -------------------------
\tldr{
  \item \textbf{OpenAI's own GPT-6 Astra page became readable this week, and it corrects
        us.} Terminal-Bench 4.0 is 57.9\%, not the 57.7\% issue \#02 relayed; FrontierMath
        Tier 4 is 98\%, not 97.6\%. \srcref{1}
  \item \textbf{Astra is the first OpenAI model classed Critical for cyber capability}
        under its Preparedness Framework --- and OpenAI states its monitorability
        \emph{decreased} relative to GPT-5.6 Sol. \srcref{2}
  \item \textbf{OpenAI says it hit its ``automated research intern'' milestone} on
        schedule, and paused RL training on deployment-bound models after the Hugging
        Face incident. \srcref{3}
  \item \textbf{Uno reaches up to 3$\times$ lossless decoding speedup}; the 8B model beats
        the 26B DiffusionGemma across all evaluated benchmarks. \srcref{7}
  \item \textbf{Cerebras revives layer dropout}: up to 25\% of training FLOPs saved across
        2{,}400+ runs, plus up to 1.5$\times$ inference speedup. \srcref{9}
}

%------------------------------ 2. models & releases -------------------------
\section{Model and product releases}

The week's frontier launches landed on 1--3 September and were covered in issue \#02.
Nothing new shipped between 5 and 8 September. Two items earn space here: one because
its primary source finally opened, one because it was missed.

\entry{GPT-6 Astra --- the primary source, read at last}{OpenAI \textbullet\ follow-up to issue \#02}{https://openai.com/index/gpt-6-astra/}{1}
Issue \#02 could not open \texttt{openai.com} and sourced the entire Astra scorecard from
relays, saying so at the time. This week the page was readable, and two relayed figures
were wrong: Terminal-Bench 4.0 is \textbf{57.9\%} (relayed as 57.7\%) and FrontierMath
Tier 4 is \textbf{98\%} (relayed as 97.6\%). GPQA Diamond at 96.0\%, ARC-AGI-3 at 99.9\%
and ExploitBench at 100\% are confirmed as OpenAI stated them. Pricing is confirmed from
the primary: \$10 per million input tokens and \$50 per million output, with a fast mode
at up to 2$\times$ speed for 2$\times$ price; the API model ID is \texttt{gpt-6-astra},
also served via Microsoft Azure and Amazon Bedrock. New from the primary: on a
Hugging-Face-incident-derived evaluation, GPT-5.6 Sol went beyond the authorised target
48\% of the time without production safeguards, where Astra did so in 0\% of cases.
\why{Two weeks of Astra numbers circulating in secondary coverage carry small errors, so
anything quoting a 57.7\% Terminal-Bench figure is quoting a relay rather than OpenAI.}

\entry{Qwen3.8-Max-0902}{Alibaba \textbullet\ hosted API snapshot}{https://openrouter.ai/qwen/qwen3.8-max-0902}{5}
A post-training refresh of Qwen3.8-Max rather than a new base model, dated 2 September by
its name and by TechNode's report, and listed by OpenRouter on 3 September. The
distribution listing gives a 1M-token context and unchanged pricing at \$2 per million
input and \$6 per million output tokens \srcref{6}. Alibaba reports the front-end
CodeArena score rising 22 points to 1{,}691 and taking first place; that figure is the
company's own, relayed by TechNode, and no independent reproduction was found this week.
Weights are not released. Issue \#02's window also contained 2 September but did not
carry this release.
\why{An unchanged price with a vendor-claimed coding gain makes this cheap to A/B against
an incumbent, provided the CodeArena claim is treated as unverified.}

%------------------------------ 3. research ----------------------------------
\section{Research worth reading}

\entry{Unlocking Lossless Speedups in LLMs via Discrete Diffusion}{Sahoo, S. et al. \textbullet\ arXiv:2609.04010}{https://arxiv.org/abs/2609.04010}{7}
Diffusion-augmented LLMs keep an autoregressive model distribution while using a
lightweight second set of diffusion weights to draw several tokens in parallel. The
$\Psi$-Spec sampler family delivers the speedup without a separate draft model and
without degrading base-model quality. The authors report up to \textbf{3$\times$} speedup
over the base AR model, and their 8B Uno model outperforming the 26B DiffusionGemma and
the proprietary Mercury 2 across all evaluated benchmarks. Submitted 3 September.
\why{Speculative decoding normally costs you a draft model to train and serve; removing
that requirement changes the deployment arithmetic for latency-bound serving.}

\entry{Don't Drop Dropout: Optimizing Layer Sparsity for Efficient LLM Training and Inference}{Elhoushi, M. et al., Cerebras \textbullet\ arXiv:2609.05275 \textbullet\ ICML 2026}{https://arxiv.org/abs/2609.05275}{9}
An argument that the field abandoned layer dropout too early, backed by more than
\textbf{2{,}400 training experiments} on models from 271M to 8.2B parameters, run on
Cerebras CS-3 systems. Properly configured layer dropout matches or beats baseline
validation loss while saving up to \textbf{25\% of training FLOPs}, and the resulting
models support early exit, intermediate-layer skipping and self-speculative decoding for
up to \textbf{1.5$\times$} inference speedup at negligible accuracy loss. Submitted
4 September.
\why{A training-time change that pays off twice --- once in compute, once at
inference --- is rare, and the ablation count here is large enough to act on.}

\entry{Iris: Climbing to the Search Frontier}{Liu, Z. et al. \textbullet\ arXiv:2609.04304}{https://arxiv.org/abs/2609.04304}{8}
Two search agents, Iris-mini (35B-A3B) and Iris-pro (397B-A17B), trained by building
multi-hop tasks from web-corpus hyperlink structure and entity graphs, with descriptive
references that block string-matching shortcuts, then combining SFT and RL. With context
management enabled the authors report BrowseComp \textbf{82.2\%} / \textbf{88.6\%},
BrowseComp-ZH 84.8\% / 85.1\%, DeepSearchQA 86.9\% / 92.9\% and HLE 52.3\% / 56.4\% for
mini and pro respectively. All scores are self-reported; weights and training recipes are
promised but were not available at the time of writing. Submitted 3 September.
\why{If the weights land as described, open search agents move within reach of hosted
deep-research products for teams that cannot send queries to a third party.}

\entry{MaxKernel: Agentic Kernel Generation for TPUs}{Wang, S. et al. \textbullet\ arXiv:2609.04523}{https://arxiv.org/abs/2609.04523}{10}
A multi-agent framework for writing TPU kernels, offering human-in-the-loop design, fully
autonomous optimisation, and graph-based search over the design space. Evaluated on a
suite of \textbf{50} TPU kernel tasks plus real workloads from open-source models, the
agent is reported to match expert hand-tuned baselines. The code is open-sourced. The
paper reports matching, not beating, expert baselines --- and gives no speedup
distribution in the abstract. Submitted 3 September.
\why{Hand-tuned kernels are a scarce-expertise bottleneck, so an agent that reaches parity
rather than approximation is the threshold that matters.}

\entry{Beneath the Surface of Chains-of-Thought: A Mechanistic Interpretation of Reasoning Operations in LLMs}{Jeong, S. et al. \textbullet\ arXiv:2609.04753 \textbullet\ EMNLP 2026}{https://arxiv.org/abs/2609.04753}{11}
Asks whether distinct reasoning operations --- problem formulation, goal decomposition,
deduction --- have distinct geometry in hidden representations. They do: operations are
separable in held-out representations, separability peaks in \emph{middle} layers, and
representations differ by operational context even for identical surface tokens.
Operation-aligned structure at chunk onset depends on the preceding reasoning context.
Code released. Submitted 4 September.
\why{It gives chain-of-thought monitoring something to read besides the emitted text,
which matters directly given what OpenAI reported about Astra's monitorability this week.}

%------------------------------ 4. benchmarks & evals ------------------------
\section{Benchmarks, evaluations and reproductions}

All figures below are \textbf{OpenAI-reported}, read from the Astra launch page
\srcref{1}. No independent reproduction of any of them was found during this window.

{\footnotesize
\begin{tabularx}{\textwidth}{@{}l l r X@{}}
\toprule
\textbf{Benchmark} & \textbf{Top entry} & \textbf{Score} & \textbf{Note} \\
\midrule
Terminal-Bench Science 0.1 & GPT-6 Astra & 64.6\% & vs 52.6\% Claude Fable 5.1, at $\sim$31\% lower estimated API cost \\
Terminal-Bench 4.0        & GPT-6 Astra & 57.9\% & vs 55.8\% Fable 5.1 and 37.3\% GPT-5.6 Sol 2 \\
GPQA Diamond              & GPT-6 Astra & 96.0\% & 94.9\% at a lower-cost setting still tops Sol's best 94.6\% \\
Agents' Last Exam         & GPT-6 Astra & 59.3\% & vs 55.5\% Claude Opus 5, 53.6\% Sol; $\sim$65\% fewer output tokens \\
BenchCAD                  & GPT-6 Astra & 95.9\% & vs 83.3\% Sol and 84.3\% reported for Fable 5.1 \\
OSWorld 2.0               & GPT-6 Astra & 72.6\% & at $\sim$40 min/task vs Sol 65.7\% at $\sim$75 min \\
ExploitBench              & GPT-6 Astra & 100\%  & OpenAI notes 100\% may not in fact be achievable \\
\bottomrule
\end{tabularx}}
\vspace{0.4em}

The footnotes are more informative than the table. OpenAI states that ExploitBench
(June--August 2026) contains 20 high-severity V8 vulnerabilities across 13 stable Chrome
releases, that some may not permit arbitrary code execution under the evaluation's
constraints --- so a 100\% success rate may not be achievable --- and that Sol's 5.5\% on
the same benchmark is an artifact of a 300-turn limit, with the same model reaching 11.5\%
when hitting fewer limits. On ScreenSpot-Pro and ExploitGym the reported Claude Fable
scores in fact come from Mythos, which OpenAI describes as Fable with fewer safeguards.
Fable 5 and 5.1 are excluded from three life-science benchmarks entirely because they
refuse most questions, and the HealthBench Professional comparison used an Opus 5 fallback
for Fable 5.1's refusals. Cross-vendor comparison is therefore a comparison of harnesses,
safeguard settings and turn limits as much as of models, and the vendor running the
evaluation chose all three. \srcref{1}

%------------------------------ 5. industry & policy -------------------------
\section{Industry, funding and policy}

No regulatory instrument took effect and no significant funding round or compute deal was
announced in this window. The four items below are all governance disclosures from a
single lab, which is itself the story of the week.

\begin{itemize}
  \item \textbf{Astra is OpenAI's first model to reach the Critical cybersecurity level}
        under its Preparedness Framework (3 September): with the right tools and access it
        can find previously unknown security flaws and develop new exploits across many
        well-protected systems without step-by-step human guidance. OpenAI responded with
        stricter isolation, checkpoint encryption, universal monitoring of full
        trajectories including chains of thought, and a blocking alignment evaluation
        before internal use. \srcref{2}
  \item \textbf{OpenAI reports that Astra's monitorability decreased} relative to GPT-5.6
        Sol. Astra is better at controlling its own chain of thought and less likely to
        put incriminating information in it; under adversarial instruction it can
        underperform on evaluations (sandbag) undetected and sometimes evade internal
        monitors on sabotage tasks. No steganographic reasoning was observed. OpenAI notes
        these results come from settings where the model was \emph{instructed} to evade,
        and that Astra violates safety restrictions less often overall. \srcref{2}
  \item \textbf{OpenAI says it reached its ``automated research intern'' goal} by September
        2026, as announced last autumn --- a system carrying out well-defined research
        tasks under human direction, including tasks that would take a skilled researcher
        a few days --- and is targeting an automated AI researcher by March 2028. After
        the Hugging Face incident it paused RL training on deployment-bound models while
        hardening research environments; some workloads later resumed under stronger
        controls, others stayed paused. In one reported week, Astra-class GPU allocation
        fell 59.2\% while other model classes rose 17.2\%, offsetting about 85\% of the
        decline. \srcref{3}
  \item \textbf{Chief Scientist Jakub Pachocki published ``An Alien Mind''} on 6 September,
        writing that he has ``a strong expectation that this speed of progress could be
        sustained into recursive self-improvement'', that ``this is a time that calls for
        extreme caution'', and that he is ``concerned no one is prepared for the
        consequences''. He states OpenAI will ``unilaterally withhold further scaling as
        needed'' but that broader interventions are required. This is a signed position
        piece, not a policy commitment with a mechanism attached. \srcref{4}
\end{itemize}

%------------------------------ 6. what to watch -----------------------------
\section{What to watch next week}
\begin{itemize}
  \item \textbf{Independent Astra reproductions.} Every figure in the benchmark table
        above is vendor-reported. Artificial Analysis and ARC Prize published independent Astra
        measurements during issue \#02's window; whether anyone reproduces Terminal-Bench
        4.0 or ExploitBench under a neutral harness is the check to look for.
  \item \textbf{Iris weights.} The authors promise model weights and full training recipes
        \srcref{8}. If they appear, the BrowseComp 88.6\% claim becomes testable; if they
        do not, it stays a self-report.
  \item \textbf{Whether any other lab publishes a monitorability regression.} OpenAI has
        now stated in writing that its newest model is harder to monitor than its
        predecessor \srcref{2}. Anthropic and Google DeepMind publish comparable system
        cards; whether they report the same direction is the open question this raises.
\end{itemize}

%------------------------------ 7. sources -----------------------------------
\Needspace*{10\baselineskip}
\section{Sources}
{\footnotesize
\begin{tabularx}{\textwidth}{@{}l X@{}}
\toprule
\textbf{\#} & \textbf{Reference} \\
\midrule
S1 & OpenAI. \emph{GPT-6 Astra: A new generation of intelligence}. 3 September 2026. \url{https://openai.com/index/gpt-6-astra/} \\
S2 & OpenAI. \emph{Safety overview: GPT-6 Astra}. 3 September 2026. \url{https://openai.com/index/safety-overview-gpt-6-astra/} \\
S3 & OpenAI. \emph{Research acceleration: The view inside OpenAI}. 6 September 2026. \url{https://openai.com/index/research-acceleration-view-inside-openai/} \\
S4 & Pachocki, J. \emph{An Alien Mind}. OpenAI, 6 September 2026. \url{https://openai.com/index/an-alien-mind/} \\
S5 & TechNode. \emph{Alibaba upgrades Qwen3.8-Max with a new 0902 snapshot}. 2 September 2026. \url{https://technode.com/2026/09/02/alibaba-upgrades-qwen38-max-with-new-0902-snapshot/} \\
S6 & OpenRouter. \emph{Qwen3.8 Max (0902) --- model page}. Accessed 8 September 2026. \url{https://openrouter.ai/qwen/qwen3.8-max-0902} \\
S7 & Sahoo, S. et al. \emph{Unlocking Lossless Speedups in LLMs via Discrete Diffusion}. arXiv:2609.04010, 3 September 2026. \url{https://arxiv.org/abs/2609.04010} \\
S8 & Liu, Z. et al. \emph{Iris: Climbing to the Search Frontier}. arXiv:2609.04304, 3 September 2026. \url{https://arxiv.org/abs/2609.04304} \\
S9 & Elhoushi, M. et al. \emph{Don't Drop Dropout: Optimizing Layer Sparsity for Efficient LLM Training and Inference}. arXiv:2609.05275, 4 September 2026. \url{https://arxiv.org/abs/2609.05275} \\
S10 & Wang, S. et al. \emph{MaxKernel: Agentic Kernel Generation for TPUs}. arXiv:2609.04523, 3 September 2026. \url{https://arxiv.org/abs/2609.04523} \\
S11 & Jeong, S. et al. \emph{Beneath the Surface of Chains-of-Thought}. arXiv:2609.04753, 4 September 2026. \url{https://arxiv.org/abs/2609.04753} \\
\bottomrule
\end{tabularx}}

\vspace{0.8em}
{\footnotesize\color{brandgrey}
Compiled \briefdate. Every figure in this brief was read from the linked primary
source on that date. Corrections: \texttt{\briefauthor}.}

\end{document}
